\paragraph{State transitions and bounded guard width.}
Let $G$ be a directed graph on a subset of the bit strings of length
$m$, with a designated start state $s_0$. As above, use two groups
of $m$ rail pairs and write $A_s,B_s$ for the coalitions selecting the
bits of $s$. The loss at $X$ is the posting $a:A_{s_0}$, and the
positive request is $X$ when the rail inputs are all $c$-atoms.

At a rail source, allow losses only if every pair in its group has a
$c$-posting. If every pair has exactly one, the other rails specify an
address $s$, which includes the head. In the $a$-group require that
$s$ is a state and all of $B_s$ is posted at $b$. In the $b$-group
require that $s$ is a state and all of $A_t$ is posted at $a$ for some
arc $s\to t$. These requirements specify exactly the losses on unique
addresses. Where a pair has two $c$-postings, any extension consistent
with antitonicity of the source predicate can be used. For instance,
one can take the upward closure of the losses specified on unique
addresses. Use the same loss predicate on all proper observer facets.

\begin{theorem}[State transitions forced by source losses]
\label{core:thm:guard-transition-depth}
These bounds and the distinguished positive are feasible exactly when
the part of $G$ reachable from $s_0$ is acyclic. If $H$ is the largest
number of vertices on a path from $s_0$ in that part, the least number
of owner slots is between $2H$ and $2H+1$.
\end{theorem}
\begin{proof}
Write $L(a,A_s)$ and $L(b,B_s)$ for the last attachments in a
linearized marked spine that win their indicated tests. As in
\cref{core:thm:compact-depth-upper-bounds}, the start loss forces
$L(a,A_{s_0})$, and the all-$c$ positive forbids a $c$-attachment.
For each state $s$, post $A_s$ at $a$ and $B_s$ at $b$, giving
the other rail inputs $\{c\}$ and giving $X$ a menu omitting $a$.
The bounds omit every $a$-source there. If $L(a,A_s)$ exists, it
therefore forces a strictly later $L(b,B_s)$. For each arc $s\to t$,
use the corresponding comparison with $B_s$ at $b$ and $A_t$ at $a$,
omitting $b$ from $X$. In each comparison use the full complement
of the omitted owner as $X$'s menu. Now $L(b,B_s)$ forces a strictly later
$L(a,A_t)$. The selected spine inherits both omissions from the
realizer and its neutral residue, just as in the previous proof.
A reachable cycle is impossible. Following a path of $H$ states
forces $2H$ distinct attachments, with one per slot on the
linearized path.

Conversely assume there is no reachable cycle. Give each reachable
state the maximum length of a path from $s_0$ to it. In increasing
rank order attach, in one $a$-fold, all the $c$-folds on the $A_s$
of that rank, and then in one $b$-fold all the $c$-folds on the
corresponding $B_s$. Begin with $X$. A rail mark at a loss can come
only from the helper of the unique $c$-free address of its group,
if that address is reachable. In the $a$-group the $b$-fold of the
same rank then absorbs it. In the $b$-group some successor's $a$-fold
absorbs it at a later rank. No subsequent helper can reintroduce that
mark. If a pair has two $c$-postings there is no helper that carries
it. The $a$-fold at the start absorbs at the required loss of $X$,
and no helper can reintroduce $X$. At the all-$c$ positive none of
the $a$- or $b$-folds can suppress the mark. All local $c$-folds fit
one final slot, after padding their branches by unary folds. The
schedule has at most $2H+1$ slots.
\end{proof}

\begin{corollary}[Exponential depth at prime width five]
\label{core:cor:width-five-depth}
For every fixed $K\ge3$ and $l\ge2$ there are feasible bounds on
$8l+1$ inputs, given by complete canonical prime lists with
$O_K(l^2)$ posting occurrences and at most five per guard, for which
one positive requires $2l\,2^l$ owner slots or more. Consequently
completion depth is exponential even at this fixed prime width.
\end{corollary}
\begin{proof}
Use states of a counter that handles one bit at each step. A state
has $l$ data bits $x_0,\ldots,x_{l-1}$, a carry bit $d$, and $l-1$
stage flags $p_1,\ldots,p_{l-1}$. Interpret $p_0=1,p_l=0$ as fixed
endpoints. Valid flags form a nonincreasing list, so the active stage
is the unique $i$ with $p_i=1,p_{i+1}=0$. At stage $i$ replace
$x_i$ by $x_i\mathbin{\mathrm{xor}}d$, leaving the other data bits
unchanged. Advance the stage, keeping carry $d\land x_i$. At the
last stage halt if $d=x_{l-1}=1$; otherwise wrap to stage zero with
carry one. Start with data zero, stage zero and carry one. Each full
sweep increments the data by one. On the last sweep the overflow
halts at its final stage, so the path from the start has $l\,2^l$
states. Its stage-zero states have distinct data, and other states
cannot repeat because a deterministic path with a repeated state
would cycle and could never halt.

There are $m=2l$ address bits. For a head's own group retain the
pair clauses for $c$ and use the $C$-literals of
\cref{core:thm:compact-depth-upper-bounds} to test its address.
Validity is checked by adjacent implications on the flags. At a
unique address the $a$-side uses equality of code bits, as before.
At a unique $b$-address, allow a transition precisely when the
state is valid and the three bits for overflow do not give a halt.
Require the corresponding next-state rails at $a$ by local clauses.
For an output data bit its old bit, the carry and the two adjacent
flags suffice. Each next stage flag tests its predecessor and the
last flag. For next carry, make a rule per stage using the two flag
tests selecting that stage, the old carry and the stage's data bit.
The wrap rule sets carry one. Expand each rule by taking the
opposite $C$-literals of its premises and adjoining the required
$a$-posting. At most four premises occur in each rule, with only
constantly many truth-table entries per bit or stage. Thus there
are $O(l)$ clauses of width at most five per rail head, and the loss
at $X$ contributes $O(l)$ singleton guards. On exact addresses the
clauses impose exactly the transitions of this counter. Doubled
$c$-postings may also satisfy the clauses, which is an allowed
extension in \cref{core:thm:guard-transition-depth}.

All literals have one fixed color at each input within each loss
predicate, so minimizing the clauses on each observer facet gives
the canonical lists, as in \cref{core:thm:compact-depth-upper-bounds}.
The state transition theorem gives the required number of slots;
\cref{core:thm:completion-fold-schedule} yields the depth bound.
\end{proof}
